\documentclass[10pt,a4paper]{amsart}
\usepackage[margin=19mm]{geometry}
\usepackage{amsmath,amssymb,mathtools}
\usepackage[scr=rsfs,cal=euler]{mathalfa}
\usepackage{xfrac}
\usepackage[unicode,hidelinks,bookmarks=false]{hyperref}
\newcommand{\ie}{i.e.\ }
\newcommand{\cf}{cf.\ }
\newcommand{\resp}{respectively\ }
\newcommand{\Subsection}{\subsection}
\providecommand{\nobreakdash}{\nobreak\hbox{-}}
\setlength{\emergencystretch}{2em}
\multlinegap0pt
\usepackage{bm}
\usepackage{bbm}
\usepackage{dsfont}
\usepackage{xspace}
\usepackage{cancel}
\usepackage{mathtools}
\usepackage{amsthm}
\makeatletter
\@ifundefined{Definition}{\newtheorem{Definition}{Definition}}{}
\@ifundefined{Corollary}{\newtheorem{Corollary}[Definition]{Corollary}}{}
\makeatother
\newcommand{\dd}{\mathrm{d}}
\newcommand{\hz}{\mathrm{HZ}}
\title[Magnetic billiards and the Hofer-Zehnder capacity of disk ta]{Magnetic billiards and the Hofer-Zehnder capacity of disk tangent bundles of lens spaces}
\author{Johanna Bimmermann, Levin Maier}
\date{Source: arXiv:2403.06761v2 (CC BY 4.0)}
\begin{document}
\maketitle

\noindent\emph{Source and licence.} Magnetic billiards and the Hofer-Zehnder capacity of disk tangent bundles of lens spaces. Version arXiv:2403.06761v2 (CC BY 4.0), distributed under the Creative Commons Attribution 4.0 International licence, as recorded in the arXiv licence metadata for this exact version and shown on the abstract page https://arxiv.org/abs/2403.06761v2. Full rights notice: licenses/arxiv-2403.06761-CC-BY-4.0.txt.\par\smallskip

\noindent\textit{Editorial scope.} Counted unit: the Corollary whose source label is \texttt{lower bound Hofer Zehnder capacity} in the article (source0.tex lines 513--517, source label \texttt{lower bound Hofer Zehnder capacity}), delivered together with the article's own complete original proof of it (source0.tex lines 518--541, one \texttt{proof} environment of 24 lines). No mathematics is written, repaired or re-derived: statement and proof are the article's own text line for line. Required context retained uncounted: shift of zero section (lines 294--296). Every definition, display and label of those lines is retained unchanged. Omitted from this package and not redistributed: 1--293, 297--512, 542--796. Nothing inside a retained excerpt is omitted or altered: every retained body is delivered exactly as the article's lines are written, so no omission receipt of any kind accompanies this package. Citations: the retained excerpts carry no citation command at all, so no bibliography entry of the article is transcribed and no reference is redistributed. Reference adaptation: none was needed and none was made; every reference inside the delivered unit resolves inside it, so no unresolved reference or citation is produced. Printed slips kept literal: no printed word, symbol, hypothesis or number of the counted statement or of its proof was altered. Layout and class: the article's own class is \texttt{article} with packages including geometry, lineno, amssymb, amsmath, amsthm, bbold, epsfig, faktor, graphicx, graphics, pifont, float, subfigure, multirow; the wrapper is the lane's pinned \texttt{amsart} editorial wrapper at 10pt on A4, which restates the theorem-like environments the retained text needs and adds the article's own macro definitions that the retained text needs (\texttt{\textbackslash dd}, \texttt{\textbackslash hz}), with extra wrapper packages (bm, bbm, dsfont, xspace, cancel, mathtools). The delivered numbering is the unit's own. Assets: no retained span contains a figure environment, a graphics-inclusion command, a PDF-page inclusion, a table or a TikZ picture; no image file is copied into this package, the package declares no asset dependency and no image file is redistributed with it. Licence: version arXiv:2403.06761v2 of this article is distributed under the Creative Commons Attribution 4.0 International licence, as recorded in the arXiv licence metadata for this exact version and shown on the abstract page https://arxiv.org/abs/2403.06761v2. A full rights notice is delivered at licenses/arxiv-2403.06761-CC-BY-4.0.txt. Characters: every retained excerpt is pure ASCII, so the delivered engine, XeLaTeX with its default Latin Modern text font, reports no missing character.
\begin{equation}\label{shift of zero section}
\iota: (D_{1-2\varepsilon}S^3, \dd\lambda-\varepsilon\pi^*\dd\alpha)\hookrightarrow (D_1 S^3,\dd\lambda):\ (x,v)\mapsto (x,v+\varepsilon ix).
\end{equation}

\begin{Corollary}\label{lower bound Hofer Zehnder capacity}
The Hofer-Zehnder capacity $c_{\hz}(D_1L(p;1),\dd\lambda)$ is bounded from below by  
$$c_{\hz}(D_1L(p;1),\dd\lambda)\geq pl=2\pi,$$
    where $l$ denotes the length of the shortest prime geodesics on $L(p;1)$.
\end{Corollary}

\begin{proof} By monotonicity of capacities and using the shift     of the zero-section \eqref{shift of zero section} we find
    $$
    c_{HZ}(D_{1-2\varepsilon}L(p;1),\dd\lambda-\varepsilon\pi^*\dd\alpha)\leq c_{HZ}(D_1L(p;1),\dd\lambda).
    $$
    Denote the sub level set of $H_\varepsilon$ by
    $$
    U_\varepsilon=\left\lbrace (x,v)\in TL(p;1)\ \Big\vert\ \ H_\varepsilon (x,v) \leq 2\pi(1-2\varepsilon)\right\rbrace.
    $$
    Observe that $U_\varepsilon\subset D_{1-2\varepsilon}L(p;1)$. Now we can modify
    $$H_\varepsilon: U_\varepsilon\to \left [\frac{\sqrt\varepsilon}{2}, 2\pi(1-2\varepsilon)\right ]$$ 
    by composing it with another function 
    $$
    f_{\varepsilon}: \left [\frac{\sqrt\varepsilon}{2}, 2\pi(1-2\varepsilon)\right ]\to [0,\infty),
    $$
    with $0\leq f_{\varepsilon}^{\prime}<1-2\mathcal{O}(\varepsilon)$, $\min(f_{\varepsilon}\circ H_\varepsilon)=0$ and $\max(f_{\varepsilon}\circ H_\varepsilon)=\max(H_{\varepsilon})-\mathcal{O}(\varepsilon)$, so that $f_{\varepsilon}\circ H_\varepsilon$ is admissible. Indeed the periods then satisfy
    $$
    T_{f_{\varepsilon}\circ H_\varepsilon}=\frac{T_{H_\varepsilon}}{f^{\prime}_{\varepsilon}(H_\varepsilon)}>1.
    $$
    In particular, we obtain
    $$
    c_{HZ}(D_1L(p;1),\dd\lambda)\geq c_{HZ}(U_\varepsilon,\dd\lambda-\varepsilon\pi^*\dd\alpha)\geq \max(f_\varepsilon\circ H_\varepsilon)= 2\pi -\mathcal{O}(\varepsilon).
    $$
    Taking $\varepsilon\to 0$ yields the corollary as the sphere covering the lens space was of radius $1$ and therefore its prime geodesic have length $2\pi$, while the shortest prime geodesics on $L(p;1)$ have length $l=2\pi/p$.
\end{proof}

\end{document}
